\documentclass[../DoAn.tex]{subfiles}
\begin{document}

\section{Thiết kế kiến trúc}
\subsection{Lựa chọn kiến trúc phần mềm}
Hệ thống KidGo được xây dựng dựa trên \textbf{Kiến trúc Nguyên khối phân tầng (Layered Monolithic Architecture)} kết hợp với mô hình \textbf{MVC biến thể dành cho RESTful API} (còn gọi là mô hình Controller-Service-Model).

\textbf{Lý thuyết kiến trúc cơ bản:}
Kiến trúc phân tầng (Layered Architecture) chia tách ứng dụng thành các lớp (layers) có vai trò và trách nhiệm riêng biệt. Các lớp ở tầng cao (như giao diện, điều khiển) sẽ gọi xuống các lớp ở tầng thấp hơn (như nghiệp vụ, dữ liệu). Việc áp dụng kiến trúc này giúp tăng tính module hóa, dễ dàng bảo trì và cô lập lỗi (Separation of Concerns). Đối với các hệ thống Web API, mô hình MVC truyền thống (Model-View-Controller) thường được tùy biến: thay vì Server trả về giao diện HTML (View) trực tiếp, Server sẽ trả về dữ liệu định dạng JSON, và phần "View" thực sự sẽ do các ứng dụng Client (React, Mobile App) đảm nhận.

\textbf{Áp dụng vào hệ thống KidGo:}
Trong dự án này, toàn bộ mã nguồn phía Backend được gói gọn trong một khối duy nhất (Monolithic) để dễ dàng triển khai và quản lý trong giai đoạn đầu phát triển. Hệ thống tuân thủ chặt chẽ mô hình Controller-Service-Model (CSM) với các thành phần cụ thể được ánh xạ như sau:

\begin{itemize}
    \item \textbf{Tầng View (Client-side):} Không nằm trực tiếp trên Backend. Vai trò hiển thị giao diện và tương tác với người dùng (Phụ huynh, Tài xế) được giao cho ứng dụng Frontend xây dựng bằng thư viện React. Giao diện này giao tiếp với Backend thông qua các RESTful API và WebSockets.
    \item \textbf{Tầng Controller (Bộ điều khiển):} Tương ứng với thư mục \texttt{controllers/} trong mã nguồn. Nhiệm vụ của tầng này là tiếp nhận các HTTP Request từ tầng Route, trích xuất dữ liệu tham số đầu vào, gọi đến tầng Service để xử lý nghiệp vụ, sau đó định dạng kết quả và trả về cho Client dưới dạng HTTP Response (JSON).
    \item \textbf{Tầng Service (Lõi nghiệp vụ):} Tương ứng với thư mục \texttt{services/}. Đây là một tầng bổ sung quan trọng, nhằm cô lập toàn bộ các quy tắc nghiệp vụ phức tạp (như tính toán chuyến đi, xác thực OTP, xử lý gọi xe). Việc tách Service giúp Controller không bị phình to logic và dễ dàng tái sử dụng mã nguồn.
    \item \textbf{Tầng Model (Dữ liệu):} Tương ứng với thư mục \texttt{models/}. Tầng này sử dụng thư viện Mongoose (ODM) để định nghĩa cấu trúc (Schema) của các thực thể (như \texttt{User}, \texttt{Trip}, \texttt{Booking}) và trực tiếp thực hiện các thao tác truy vấn vào cơ sở dữ liệu MongoDB.
    \item \textbf{Tầng Middleware và Sockets:} Hệ thống bổ sung thêm các tầng cắt ngang tương ứng với thư mục \texttt{middlewares/} (để xử lý xác thực token, kiểm tra quyền truy cập) và \texttt{sockets/} (để duy trì luồng kết nối định vị theo thời gian thực độc lập với luồng HTTP).
\end{itemize}

Việc áp dụng kiến trúc này giúp mã nguồn của KidGo được tổ chức vô cùng mạch lạc. Các luồng dữ liệu đi một chiều từ Route $\rightarrow$ Controller $\rightarrow$ Service $\rightarrow$ Model, không bị chồng chéo, giúp hệ thống dễ dàng bảo trì hoặc mở rộng chức năng trong tương lai.

\subsection{Thiết kế tổng quan}
Sinh viên vẽ biểu đồ gói UML (UML package diagram), nêu rõ sự phụ thuộc giữa các gói (package). SV cần vẽ các gói sao cho chúng được phân theo các tầng rõ ràng, không được sắp đặt package lộn xộn trong hình vẽ. Sinh viên chú ý các quy tắc thiết kế (Các gói không phụ thuộc lẫn nhau, gói tầng dưới không phụ thuộc gói tầng trên, không phụ thuộc bỏ qua tầng, v.v.) và cần giải thích sơ lược về mục đích/nhiệm vụ của từng package. SV tham khảo ví dụ minh họa trong Hình \ref{fig:Fig1}
\begin{figure}[H]
    \centering
    \includegraphics{Hinhve/Picture1.png}
    \caption{Ví dụ biểu đồ phụ thuộc gói}
    \label{fig:Fig1}
\end{figure}
\subsection{Thiết kế chi tiết gói}
Sinh viên thiết kế và lần lượt vẽ biểu đồ thiết kế cho từng package, hoặc một nhóm các package liên quan để giải quyết một vấn đề gì đó. Khi vẽ thiết kế gói, sinh viên chỉ cần đưa tên lớp, không cần chỉ ra các thành viên phương thức và thuộc tính. SV tham khảo ví dụ minh họa trong Hình \ref{fig:Fig2}.

Sinh viên cần vẽ rõ ràng quan hệ giữa các lớp trong biểu đồ. Các quan hệ bao gồm: phụ thuộc (dependency), kết hợp (association), kết tập (aggregation), hợp thành (composition), kế thừa (inheritance), và thực thi (implementation). Các quan hệ này đều đã được minh họa trong \ref{fig:Fig2}.

Sau khi vẽ hình minh họa, sinh viên cần giải thích ngắn gọn về thiết kế của mình.

\begin{figure}[H]
    \centering
    \includegraphics{Hinhve/Picture2.png}
    \caption{Ví dụ thiết kế gói}
    \label{fig:Fig2}
\end{figure}

\section{Thiết kế chi tiết}
\subsection{Thiết kế giao diện}
Phần này có độ dài từ hai đến ba trang. Sinh viên đặc tả thông tin về màn hình mà ứng dụng của mình hướng tới, bao gồm độ phân giải màn hình, kích thước màn hình, số lượng màu sắc hỗ trợ, v.v. Tiếp đến, sinh viên đưa ra các thống nhất/chuẩn hóa của mình khi thiết kế giao diện như thiết kế nút, điều khiển, vị trí hiển thị thông điệp phản hồi, phối màu, v.v. Sau cùng sinh viên đưa ra một số hình ảnh minh họa thiết kế giao diện cho các chức năng quan trọng nhất. Lưu ý, sinh viên không nhầm lẫn giao diện thiết kế với giao diện của sản phẩm sau cùng.
\subsection{Thiết kế lớp}
Phần này có độ dài từ ba đến bốn trang. Sinh viên trình bày thiết kế chi tiết các thuộc tính và phương thức cho một số lớp chủ đạo/quan trọng nhất của ứng dụng (từ 2-4 lớp). Thiết kế chi tiết cho các lớp khác, nếu muốn trình bày, sinh viên đưa vào phần phụ lục.

Để minh họa thiết kế lớp, sinh viên thiết kế luồng truyền thông điệp giữa các đối tượng tham gia cho 2 đến 3 use case quan trọng nào đó bằng biểu đồ trình tự (hoặc biểu đồ giao tiếp).
\subsection{Thiết kế cơ sở dữ liệu}
Phần này có độ dài từ hai đến bốn trang. Sinh viên thiết kế, vẽ và giải thích biểu đồ thực thể liên kết (E-R diagram). Từ đó, sinh viên thiết kế cơ sở dữ liệu tùy theo hệ quản trị cơ sở dữ liệu mà mình sử dụng (SQL, NoSQL, Firebase, v.v.)

\section{Xây dựng ứng dụng}
\subsection{Thư viện và công cụ sử dụng}
Sinh viên liệt kê các công cụ, ngôn ngữ lập trình, API, thư viện, IDE, công cụ kiểm thử, v.v. mà mình sử dụng để phát triển ứng dụng. Mỗi công cụ phải được chỉ rõ phiên bản sử dụng. SV nên kẻ bảng mô tả tương tự như Bảng \ref{table:my_label}. Nếu có nhiều nội dung trình bày, sinh viên cần xoay ngang bảng.

\begin{table}[H]
    \centering{}
    \begin{tabular}{lll}
        \hline
        \textbf{Mục đích} & \textbf{Công cụ}       & \textbf{Địa chỉ URL}    \\ \hline
        IDE lập trình     & Eclipse Oxygen a64 bit & http://www.eclipse.org/ \\ \hline
        v.v.              & v.v.                   & v.v.                    \\ \hline
    \end{tabular}
    \caption{Danh sách thư viện và công cụ sử dụng}
    \label{fig:my_label}
\end{table}

\subsection{Kết quả đạt được}
Sinh viên trước tiên mô tả kết quả đạt được của mình là gì, ví dụ như các sản phẩm được đóng gói là gì, bao gồm những thành phần nào, ý nghĩa, vai trò?

Sinh viên cần thống kê các thông tin về ứng dụng của mình như: số dòng code, số lớp, số gói, dung lượng toàn bộ mã nguồn, dung lượng của từng sản phẩm đóng gói, v.v. Tương tự như phần liệt kê về công cụ sử dụng, sinh viên cũng nên dùng bảng để mô tả phần thông tin thống kê này.

\subsection{Minh họa các chức năng chính}
Sinh viên lựa chọn và đưa ra màn hình cho các chức năng chính, quan trọng, và thú vị nhất. Mỗi giao diện cần phải có lời giải thích ngắn gọn. Khi giải thích, sinh viên có thể kết hợp với các chú thích ở trong hình ảnh giao diện.

\section{Kiểm thử}
Phần này có độ dài từ hai đến ba trang. Sinh viên thiết kế các trường hợp kiểm thử cho hai đến ba chức năng quan trọng nhất. Sinh viên cần chỉ rõ các kỹ thuật kiểm thử đã sử dụng. Chi tiết các trường hợp kiểm thử khác, nếu muốn trình bày, sinh viên đưa vào phần phụ lục.
Sinh viên sau cùng tổng kết về số lượng các trường hợp kiểm thử và kết quả kiểm thử. Sinh viên cần phân tích lý do nếu kết quả kiểm thử không đạt.
\section{Triển khai}
Sinh viên trình bày mô hình và/hoặc cách thức triển khai thử nghiệm/thực tế. Ứng dụng của sinh viên được triển khai trên server/thiết bị gì, cấu hình như thế nào. Kết quả triển khai thử nghiệm nếu có (số lượng người dùng, số lượng truy cập, thời gian phản hồi, phản hồi người dùng, khả năng chịu tải, các thống kê, v.v.)

\end{document}
